\section{Numerical algorithm and what a trajectory means}
\label{sec:algorithm}

\subsection{Coupling scale and the implemented objective}

All evaluated instances have zero local fields. Let $K^{\rm raw}$ be the
symmetric, zero-diagonal Ising coupling matrix. For a graph with edges
listed once, the implementation defines
\begin{equation}
 \begin{aligned}
 c_K&=\sqrt{\frac{1}{N}\sum_{i,j}(K^{\rm raw}_{ij})^2}\\
    &=\sqrt{\frac{2}{N}\sum_{\{i,j\}\in E}w_{ij}^2},
 \qquad K=K^{\rm raw}/c_K.
 \end{aligned}
 \label{eq:rmscoupling}
\end{equation}
It optimizes Eq.~\eqref{eq:qefemF} with $K$, then scores binary spins
against $K^{\rm raw}$. The numerical values of $T$ and $\Gamma$ must
therefore be read in units of the normalized interaction scale. If one
multiplies the normalized free energy by $c_K$, the equivalent physical
temperature and transverse field are $c_KT$ and $c_K\Gamma$. Coupling
normalization is more than a harmless change of printed units if the
temperature and driver schedules are left fixed while $c_K$ changes.

The selected scale is 44.71018 for K2000, 0.316598 for the saved Wishart
instances, and 3.162278 for the saved Chook instances. The values and
their source configurations are retained with the result folders. G-set
uses graph-specific normalization, and the later spectral schedule also
uses the lowest eigenvalue of the normalized graph matrix.

\subsection{Initialization, schedules, and optimization}

A trial starts with $B$ independent replicas. For replica $b$ and site $i$,
the two local fields are initialized as
\begin{equation}
 a^{(0)}_{b,i,d}=h_{\rm init}\,\xi_{b,i,d},
 \qquad \xi_{b,i,d}\sim\mathcal N(0,1),\quad d\in\{x,z\}.
 \label{eq:init}
\end{equation}
An isolated CPU random generator produces the stored initialization for
each seed before it is transferred to the compute device. This improves
replayability; it does not make distinct devices bitwise equivalent.
The K2000 $h_{\rm init}$ is 0.1. A trial is one full population of $B$
replicas, not one replica.

For $L$ optimizer updates, define $\tau_t=t/(L-1)$ for
$t=0,\ldots,L-1$. The simple frozen schedules have
\begin{align}
 T_t&=T_{\min}+(T_{\max}-T_{\min})(1-\tau_t),
 \label{eq:tempcurve}\\
 \Gamma_t&=\Gamma_{\min}+
 (\Gamma_{\max}-\Gamma_{\min})
 \left[1-\min\!\left(\frac{\tau_t}{f_\Gamma},1\right)\right].
 \label{eq:gammacurve}
\end{align}
The code passes the reciprocal $\beta_t=1/T_t$ to the free-energy
kernel. For K2000 and Wishart, $f_\Gamma=0.5$, so the transverse field
vanishes halfway and the remaining updates optimize with $\Gamma=0$.
For Chook, $f_\Gamma=1$. The G-set runs also use quadratic cooling or the
spectral construction of Eq.~\eqref{eq:spectralinvert}, according to their
graph-level plans. A schedule is applied to the objective at each update;
it is not a simulated measurement history or Schr\"odinger evolution.

For the frozen K2000, Wishart, and Chook runs, automatic differentiation
computes Eq.~\eqref{eq:fieldgrad}, and PyTorch RMSprop updates all field
coordinates. The recorded implementation uses ordinary \emph{coupled}
weight decay, so its elementwise updates can be written as
\begin{align}
 \widetilde g_t&=\nabla_{\bm a}\mathcal F_t+\lambda_{\rm wd}\bm a_t,
 &v_t&=\alpha v_{t-1}+(1-\alpha)\widetilde g_t^{\odot2},
 \label{eq:rmsprop1}\\
 b_t&=\mu b_{t-1}+\frac{\widetilde g_t}{\sqrt{v_t}+\eps_{\rm opt}},
 &\bm a_{t+1}&=\bm a_t-\eta b_t.
 \label{eq:rmsprop2}
\end{align}
These equations describe the non-centered, non-Nesterov mode used in the
saved plan. The operations are elementwise in the optimizer state, except
for the matrix product in the gradient. Here
$\eps_{\rm opt}=10^{-8}$. Weight decay changes the numerical trajectory
and thus must be reported with the learning rate and momentum.
Because each replica has separate coordinates, the batch is a parallel
collection of independent optimizations under the same graph and schedule.
The implementation backpropagates a sum of per-replica free energies; it
does not divide the gradient by $B$.

For clarity, the complete frozen-run logic is:
\begin{enumerate}[label=\arabic*.,leftmargin=2em]
 \item Load and validate the graph or planted Ising instance; form normalized
       couplings $K=K^{\rm raw}/c_K$.
 \item For each declared seed, construct $B\times N\times2$ Gaussian
       initial fields and arrays $\{T_t,\Gamma_t\}_{t=0}^{L-1}$.
 \item For every update, map fields through Eq.~\eqref{eq:fieldmap},
       evaluate Eq.~\eqref{eq:qefemF}, and apply Eq.~\eqref{eq:rmsprop1}
       --\eqref{eq:rmsprop2} or the declared alternative optimizer.
 \item After exactly $L$ updates, decode all $B$ final states and rescore
       each binary vector with the raw instance objective.
 \item Record the best final replica for that seed; aggregate these trial
       endpoints using the declared mean, best, gap, and exact-hit metrics.
\end{enumerate}
The K2000 population experiment uses nested prefixes of the same initialized
8,192-replica population for each seed. This controls the random population
when $B$ changes: a 128-replica run uses exactly the first 128 initial
states of that seed's 8,192-replica run, and similarly for larger sizes.
It does not imply that their later trajectories remain prefixes, since
floating-point batching and optimizer execution can differ.

\subsection{Binary readout and scoring}

The implemented final readout is
\begin{equation}
 \widehat s_{b,i}=\begin{cases}+1,&m^{(L)}_{z,b,i}>0,\\-1,&m^{(L)}_{z,b,i}\leq0.
 \end{cases}
 \label{eq:readout}
\end{equation}
This is equivalent to rounding the probability
$p_{b,i}=(1+m_{z,b,i})/2$ as used by the solver, including the tie at
$p=1/2$. It yields a feasible binary vector for every replica. For
MaxCut, we evaluate Eq.~\eqref{eq:cut}; for planted Ising instances we
evaluate $E_I(\widehat{\bm s}_b)$ and compare it with the supplied planted
energy. No local improvement, bit flipping, coordinate splicing, or
selection of an intermediate iterate is part of the frozen endpoint
protocol.

This distinction is crucial for interpreting trajectory figures. The
stored K2000 trace first selects the final best replica of the best
endpoint trial, then replays that \emph{same replica} from its initial
field to the final update. The curve is not the maximum cut over the
population at every step. Its temporary decreases are possible and do
not contradict final best-of-population selection. The entropy trace is
likewise the sum of site entropies for that coherent replica, divided by
$N$ when plotted.

\subsection{Arithmetic and storage}

For a dense graph, one longitudinal field evaluation is a batched
matrix product $\bm m_zK$, costing $O(BN^2)$ arithmetic, with
$O(N^2+BN)$ persistent storage before optimizer buffers and temporary
autograd tensors. An edge-list or sparse matrix contraction can reduce
arithmetic to $O(B|E|+BN)$ and graph storage to $O(|E|)$, subject to device
support and implementation details. K2000 is complete, so $|E|=N(N-1)/2$
and the dense formulation is natural. Large replica counts improve the
number of basins explored but increase memory and arithmetic nearly
linearly in $B$. Hence solution quality at $B=8{,}192$ cannot be read as
the quality of a single cheap trajectory, and unmatched wall-clock
measurements across methods are not evidence of a speedup.
